\begin{textsample}{POS Dim 1 – vem\_ed\_16 – Score 25.00 – vem\_ed\_16/t070.txt}  \label{ex:f1_pos_088}
VEm Virtual Exchange Medium Como vê as iniciativas de cooperação Sul-Sul nos Intercâmbios Virtuais? \textbf{É} explícita, nas instituições latino-americanas, por exemplo, a resistência a cooperação com o Sul Global.
Uma breve análise do \textbf{número} de \textbf{acordos} de nossas IES com \textbf{cada} região do planeta poderá, facilmente, mostrar \textbf{que} os países do Norte, com suas instituições centenárias e suas ofertas tentadoras, estão na dianteira do interesse de nossos estudantes e pesquisadores como destino de \textbf{intercâmbio}.
A possibilidade dos Intercâmbios Virtuais entre instituições do Sul abre, pois, uma via a \textbf{ser} explorada para uma quebra inicial – \textbf{que} poderá \textbf{ser} \textbf{seguida} de um desejável equilíbrio – da primazia do Norte em nossas relações bilaterais.
As \textbf{trocas} virtuais \textbf{permitirão} o \textbf{desenvolvimento} de algo \textbf{que} \textbf{é} fundamental para a cooperação Sul-Sul, a \textbf{saber}, o \textbf{conhecimento} \textbf{mútuo} dos parceiros potenciais, de sua atuação na pesquisa, de sua cultura.
Somente esse \textbf{conhecimento} \textbf{permitirá} a \textbf{construção} de sólidas parcerias entre instituições do Sul, talvez com uma reciprocidade maior \textbf{que} aquela \textbf{que} temos \textbf{observado} nas relações com o Norte.
Comente a conquista do selo FAUBAI-BRaVE pelo CPS, em outubro de 2022.
O selo FAUBAI-BRaVE \textbf{foi} \textbf{criado} para atestar a participação das IES membros da FAUBAI nas nossas promoções de Intercâmbio Virtual nos moldes do \textbf{Programa} FAUBAI-BRaVE.
A atribuição do selo ao Centro Paula Souza, além de \textbf{reconhecer} a qualidade de seus projetos de Intercâmbio Virtual, \textbf{representa} também o reconhecimento da FAUBAI pela participação do CPS, juntamente com a UNESP e a UFPE, na experiência piloto \textbf{que} \textbf{construiu} esse \textbf{programa}.
Como as demais IES portadoras desse selo, o CPS poderá utilizá-lo para comprovar sua participação no \textbf{programa} diante dos parceiros \textbf{que} vier a conquistar e, também, terá direito a condições especiais nas promoções \textbf{que} o \textbf{programa} FAUBAI-BRaVE vier a lançar no \textbf{futuro}, como ofertas de parcerias estabelecidas pela FAUBAI, gratuidades ou descontos em formações \textbf{específicas} para a área, entre outros.
Como vislumbra o \textbf{futuro} da internacionalização?
O cenário macropolítico internacional aparenta um recrudescimento de ideias totalitárias e, por vezes, nacionalistas, contra o \textbf{ideal} democrático \textbf{que} faz transpor fronteiras.
Essa impressão acentuou-se no período da pandemia, quando as \textbf{trocas}, num primeiro \textbf{momento}, pareciam \textbf{ser} impossíveis.
Entretanto, podemos dizer sem exagero \textbf{que} o braço tecnológico da globalização se \textbf{estabeleceu} solidamente em \textbf{todas} as culturas, permitindo-nos, inclusive, enfrentar com menos sofrimento as agruras das quarentenas.
Espantosamente, as tecnologias \textbf{que}, pelo alto custo do acesso à Web e a equipamentos, sempre \textbf{foram} excludentes, tornaram-se um \textbf{fator} de inclusão de primeira ordem no mundo da educação, para nos restringirmos ao nosso campo.
Isso mostra \textbf{que} detemos um \textbf{conhecimento} capaz de nos fazer superar qualquer obstáculo, incluindo o negacionismo e o obscurantismo \textbf{que} compõem as receitas totalitárias.
Vejo, pois, o futuro da educação \textbf{superior} internacionalizada sempre mais adepto das \textbf{trocas} virtuais, \textbf{que} já provaram \textbf{ser} de grande utilidade, mesmo diante do mau uso \textbf{que} pode, por vezes, permeá-la.
Ter o mundo em nossa \textbf{casa}, com o recurso quase mágico das novas tecnologias, \textbf{é} uma oportunidade \textbf{que} não pode \textbf{ser} ignorada, pois, além dos contatos rápidos e constantes, oferece também uma via mais segura e econômica dada a \textbf{necessidade} de redução da emissão de carbono.
Não se trata de imaginar o \textbf{futuro} da internacionalização apenas em um metaverso crescente, mas sim de integrar essas possibilidades virtuais ao \textbf{processo} educacional e à \textbf{produção} científica.

% matched lemmas: acordo, cada, casa, conhecimento, construir, construção, criar, desenvolvimento, específico, estabelecer, fator, futuro, ideal, intercâmbio, momento, mútuo, necessidade, número, observar, permitir, processo, produção, programa, que, reconhecer, representar, saber, seguir, ser, superior, todo, troca
\end{textsample}
